\documentclass[10pt,a4paper]{amsart}
\usepackage[margin=19mm]{geometry}
\usepackage{amsmath,amssymb,mathtools}
\usepackage[scr=rsfs,cal=euler]{mathalfa}
\usepackage{xfrac}
\usepackage[unicode,hidelinks,bookmarks=false]{hyperref}
\newcommand{\ie}{i.e.\ }
\newcommand{\cf}{cf.\ }
\newcommand{\resp}{respectively\ }
\newcommand{\Subsection}{\subsection}
\providecommand{\nobreakdash}{\nobreak\hbox{-}}
\setlength{\emergencystretch}{2em}
\multlinegap0pt
\usepackage{bbm}
\newcommand{\on}{\operatorname}
\newcommand{\IR}{\mathbb{R}}
\newcommand{\IQ}{\mathbb{Q}}
\newcommand{\IN}{\mathbb{N}}
\newcommand{\IZ}{\mathbb{Z}}
\newcommand{\IC}{\mathbb{C}}
\newcommand{\IK}{\mathbb{K}}
\newcommand{\LC}{\mathcal{L}}
\newcommand{\I}{\mathcal{I}}
\newcommand{\J}{\mathcal{J}}
\newcommand{\Ma}{\mathbb{M}}
\newcommand{\mc}{\mathcal}
\newcommand{\mf}{\mathfrak}
\newcommand{\Fin}{\mathrm{Fin}}
\newcommand{\Exh}{\mathrm{Exh}}
\newcommand{\supp}{\mathrm{supp}}
\newcommand{\CR}{\mathcal{CR}}
\newcommand{\w}{\mathbbm}
\newtheorem{theorem}{Theorem}[section]
\newtheorem{definition}[theorem]{Definition}
\newtheorem{lemma}[theorem]{Lemma}
\newtheorem{example}[theorem]{Example}
\newtheorem{proposition}[theorem]{Proposition}
\newtheorem{remark}[theorem]{Remark}
\newtheorem{corollary}[theorem]{Corollary}
\newtheorem{question}[theorem]{Question}
\title[Zoo of ideal Schauder bases]{Zoo of ideal Schauder bases: distinguishing ideals by simple ideal Schauder bases (Theorem 4.1)}
\author{Adam Kwela and Jaros{\l}aw Swaczyna}
\date{Source: arXiv:2508.21235v1 (CC BY 4.0)}
\begin{document}
\begin{abstract}
We investigate the notion of filter (equivalently: ideal) Schauder basis of a Banach space. We do so by providing bunch of new examples of such bases that are not the standard ones, especially within classical Banach spaces ($\ell_p$, $c_0$, James' space). Those examples lead to distinguishing and characterizing filters (equivalently: ideals) in terms of Schauder bases. We investigate the relationship between possibly basic sequences and filters (equivalently: ideals) on the set of natural numbers. 
\end{abstract}
\maketitle
\noindent\emph{Source and licence.} Zoo of ideal Schauder bases. Version arXiv:2508.21235v1 (CC BY 4.0), distributed under the Creative Commons Attribution 4.0 International licence, http://creativecommons.org/licenses/by/4.0/, as recorded in the arXiv OAI licence metadata for this exact version and shown on the abstract page https://arxiv.org/abs/2508.21235v1. Full rights notice: licenses/arxiv-2508.21235-CC-BY-4.0.txt.\par\smallskip
\noindent\textit{Editorial scope.} Counted unit: the paper's distinguishing theorem, source label thm-distinguishing (source0.tex lines 510--519), printed as Theorem 4.1 in the published version and as Theorem 4.1 here, delivered together with its complete original author proof at source lines 521--566 as the single primary proof body. No mathematics is written, repaired or re-derived: the statement and the proof are the paper's own text line for line, including its own typographic forms, its own \{\bf F\}-style bold symbols, its mathtools cases* displays, its interval notation and its printed slips. The paper is already an amsart article (source0.tex line 2) with no class options, so no publisher class is replaced or shipped: the wrapper is the lane's pinned amsart editorial wrapper, and the paper's own macro block (source lines 25--42) and its own \textbackslash{}newtheorem declarations (source lines 44--51, all sharing the section-numbered theorem counter) are copied verbatim into the wrapper preamble, so every printed number of the delivered unit is produced by the paper's own counters and coincides with the published version: Sections 1 Preliminaries, 2 Ideal Schauder bases, 3 Some lemmas, 4 Banach spaces with a Schauder basis, all the theorem-like items of those sections (the two folklore lemmas, the example listing well-known ideals, the two general lemmas with their remark, the example of a normed-basis reduction, the proposition that the critical ideal of a simple basis is a P-ideal, the reduction theorem for simple bases over a normed basis, the coefficient lemmas and their corollaries) and the counted Theorem 4.1. Required context retained uncounted: the whole of the paper's Sections 1--3 (source lines 79--506) together with the heading of Section 4 (source lines 508--509). Section 1 Preliminaries supplies the notions of an ideal, of a tall ideal, of a P-ideal, of a summable ideal, of a submeasure and of a lower semicontinuous or non-pathological submeasure, the definition of ideal convergence of sequences and of series (the notation used in the counted statement and proof), the definition of a Schauder basis as a sequence of pairs with its coordinate functionals, of an unconditional basis and of the standard basis of the space c-nought, and the isomorphism lemma and the monotonicity lemma with their proofs. Section 2 Ideal Schauder bases supplies the definition of an I-Schauder basis, the critical ideal of a basis, the definition of a basis simple over another one (source lines 331--342, Definition def-simple) on which the counted statement and proof run, and the reduction theorem with its complete proof. Section 3 Some lemmas supplies the two general lemmas on the critical ideal of a simple basis and the proposition that it is a P-ideal, each with its complete proof. Every definition, lemma, proposition, theorem, corollary, example, remark, numbered display, footnote and citation of those source lines is retained. The paper's abstract (source line 74) is reproduced verbatim in the wrapper front matter and the paper's own title and author names are reproduced in the wrapper title block. Not part of this unit and not redistributed: the paper's address, email, url, thanks, keywords and subject-class front matter (source lines 56--66); the two corollaries deduced from the counted theorem and the remark that follows them (source lines 568--591), which the retained spans stop short of; Sections 5--11 (source lines 593--1415), that is the classification of ideal Schauder bases in the spaces ell-p, in c-nought and in James' space, the construction for every non-pathological analytic P-ideal, the section on avoiding the space of finitely supported sequences, the ideal version of FDD, the section on normed vector spaces and the closing list of questions; and the paper's bibliography listing, of which only the cited entries are transcribed. The delivered text cites fifteen keys (AK, CGK, Debs, RKS, Farah, Fast, Fridy, KS, Kochanek, KSW, Mazur, Salat, Solecki, SoleckiExh and Steinhaus, twenty citation commands in all), and those fifteen entries are transcribed verbatim from the paper's own in-source thebibliography (source lines 1419--1478) with the paper's printed bibliography numbers preserved (AK prints 1, CGK 5, Debs 6, RKS 7, Farah 8, Fast 9, Fridy 10, KS 14, Kochanek 15, KSW 16, Mazur 17, Salat 18, Solecki 19, SoleckiExh 20 and Steinhaus 21), so every printed citation number of the delivered unit is the published one, including the four-key citation at source line 238 whose keys print 5, 7, 14 and 15 in the paper's own order. Reference adaptation: one key, two occurrences, method reference\_only\_adaptation recorded in delivery-evidence.json. The paper refers twice to its own Section 10 Normed vector spaces, whose label S:normed lies outside this unit (source line 1321), once at source line 316 and once at source line 413; both occurrences are delivered as \textbackslash{}href links to the abstract page of this exact version with the printed number retained, so no unresolved reference is produced. Every other label and reference of the delivered spans is the paper's own and resolves inside the unit (lem-isomorphisms, lem-monotonicznosc, L:wspolczynniki, P:wspolczynniki, def-simple, general-lem1, remark-general-lem1, general-lem2, general-P-ideal and the counted thm-distinguishing), and no label occurs twice. Printed slips kept literal: no printed word, symbol, hypothesis or number of the counted statement or of its proof was altered. The paper's own slips are delivered as printed, among them the misprint Godefory for Godefroy in the entry AK, the doubled subscript in the definition of the coordinate functionals at source line 206, the phrase convergent to X at source line 198 and the paper's inconsistent spelling of base and basis. Unicode: the source contains six non-ASCII characters in four distinct code points (one right single quotation mark at source line 1039, one e-acute in the entry ARTT at source line 1423, three a-acute in the entry MezaGuzman at source line 1447 and one en dash in the entry Szar at source line 1475). None of them lies inside the retained spans or inside the fifteen transcribed entries; the retained spans and the delivered bibliography are pure ASCII, so no character of the source had to be replaced, dropped or re-encoded, no fidelity receipt for a character is needed, and the xelatex receipt reports no missing character. The paper's loads of graphicx, latexsym, accents, mathrsfs, euscript, mdwlist, todonotes, enumerate, enumitem, a4wide, physics, color, hyperref and amsthm are either already supplied by the wrapper's pinned preamble or unused by the retained spans, which contain no figure, no table, no \textbackslash{}includegraphics, no \textbackslash{}input and no custom list environment; the bbm package is loaded because the paper's own macro \textbackslash{}w is defined as \textbackslash{}mathbbm, and the mathtools package is present for the cases-star displays of the counted proof. No mathematics is written, repaired or re-derived; all mathematics of the unit is native text with no image dependency.
\section{Preliminaries}
We aim to keep our paper readable for both set theorists and functional analysts; therefore, we recall some of the basic definitions from both fields.

Throughout the paper, we use standard set-theoretic notation and identify each natural number $n\in\omega=\{0,1,2,\ldots\}$ with the set $\{0,1,\ldots,n-1\}$.

\subsection{Ideals}
We will be concerned with ideals $\I$ on the set of natural numbers $\omega$. By an ideal on a countable set $M$ we mean a family $\I \subseteq \mathcal{P}(M)$ such that:
\begin{itemize}
    \item $\I$ contains all finite subsets of $M$;
    \item if $A,B \in \I$, then also $A\cup B \in \I$;
    \item if $A \in \I$ and $B \subseteq A$, then also $B \in I$.
\end{itemize}
We will focus primarily on non-trivial ideals, that is, ideals $\I$ such that $\I \neq \mathcal{P}(M)$, and our ideals will typically be defined on $M=\omega$. 

\begin{remark}
Clearly, the notion of ideal is dual to the notion of filter: $\I$ is an ideal if and only if the family of complements of elements of $\I$ forms a filter. Therefore, choice between working in terms of ideal bases and filter bases is purely aesthetic.
\end{remark}    

We treat the power set $\mathcal{P}(M)$ as the space $2^M$ of all functions $f:M\rightarrow 2$ (equipped with the product topology, where each space $2=\left\{0,1\right\}$ carries the discrete topology) by identifying subsets of $M$ with their characteristic functions. All topological and descriptive notions in the context of ideals will refer to this topology. 

An ideal $\I$ on $M$ is tall if for any infinite $A\subseteq M$ there is an infinite $B\subseteq A$ such that $B\in\I$. Equivalently, $\I$ is tall if $\I\restriction A\neq[A]^{<\omega}$, for any $A\notin\I$, where
\[
\I\restriction A=\left\{B\cap A:B\in\I\right\}.
\]
We say that an ideal $\I$ on $M$ is generated by a family $\mathcal{A}\subseteq\mathcal{P}(M)$ if 
\[
\I=\{B\subseteq M:\exists_{n\in\omega}\ \exists_{A_0,\ldots,A_n\in\mathcal{A}}\ B\subseteq A_0\cup\ldots\cup A_n\}.
\]
An ideal $\I$ is called a P-ideal if for each sequence $(A_n)\subseteq  \I$ there is a set $A\in\I$ such that $A_n\setminus A$ is finite for all $n\in\omega$. By a result of Solecki, every analytic P-ideal is $\bf{F_{\sigma\delta}}$ (see \cite{Solecki}).

If $\{X_m: m\in M\}$ is a family of sets, then $\sum_{m\in M}X_m=\{(m,x): m\in M,x\in X_m\}$ is its disjoint sum. The vertical section of a set $A\subseteq \sum_{m\in M}X_m$ at a point $m\in M$ is defined
by $A_{(m)} = \{x\in X_m: (m,x)\in A\}$. Let $\I$ and $\J$ be ideals on $M$ and $N$, respectively, and $\I_m$, for each $m\in M$, be an ideal on $X_m$. We define the following new ideals:
\begin{itemize}
	\item $\sum_{m\in M}\I_m=\{A\subseteq \sum_{m\in M}X_m: \forall_{m\in M}\ A_{(m)}\in\I_m\}$;
 	%\item $\sum^J_{t\in T}\I_t=\{M\subseteq \sum_{t\in T}X_t:\ \{t\in T:\ M_{(t)}\notin\I_t\}\in\J\}$;
    %\item $\I\otimes\J=\sum^\I_{s\in S}\J=\{M\subseteq S\times T:\ \{s\in S:\ M_{(s)}\notin\J\}\in\I\}$;
	%\item $\I\otimes \{\emptyset\} = \{M\subseteq S\times \omega:\ \{s\in S:\ M_{(s)}\ne\emptyset\}\in \I\}$;
	\item $\{\emptyset\} \otimes  \I= \{A\subseteq \omega\times M: A_{(n)} \in \I\text{ for all }n\in \omega\}$;
    \item $\I\oplus\J=\{A\subseteq(\{0\}\times M)\cup(\{1\}\times N): A_{(0)}\in\I\text{ and }A_{(1)}\in\J\}$.
\end{itemize}

Two ideals $\I$ and $\J$ on $M$ and $N$, respectively, are isomorphic ($\I\cong\J$) if there is $f:N\to M$ such that:
\[
f^{-1}[A]\in\J\ \Leftrightarrow\ A\in\I,
\]
for all $A\subseteq M$. Two isomorphic ideals have the same descriptive complexity (see for instance \cite[Lemma 7.2]{Debs}). 

\begin{lemma}[Folklore]
\label{lem-isomorphisms}
Let $\I$ be an ideal.
\begin{itemize}
    \item[(a)] If $\I$ is not tall, then $\I\cong\I\oplus\Fin$.
    \item[(b)] If $\I\neq\Fin$, then $\I\cong\I\oplus\mathcal{P}(\omega)$.
\end{itemize}
\end{lemma}

\begin{proof}
Without loss of generality, we may assume that $\I$ is an ideal on $\omega$. 

Since $\I$ is not tall ($\I\neq\Fin$, respectively), there is $A\notin\I$ such that $\I\restriction A=\Fin\restriction A$ (an infinite $A\in\I$, respectively). Let $g:((\{0\}\times A)\cup(\{1\}\times\omega))\to A$ be any bijection. Then $f:2\times\omega\to\omega$ given by:
\[
f(i,n)=\begin{cases}
    g(i,n), & \text{if }(i,n)\in(\{0\}\times A)\cup(\{1\}\times\omega),\cr
    n, & \text{otherwise,}
\end{cases}
\]
for all $(i,n)\in 2\times\omega$, witnesses $\I\cong\I\oplus\Fin$ ($\I\cong\I\oplus\mathcal{P}(\omega)$, respectively).
\end{proof}

\begin{example}
The following is a list of some well-known ideals:
\begin{itemize}
    \item The smallest ideal is $\Fin=[\omega]^{<\omega}$.  It is an $\bf{F_\sigma}$ non-tall P-ideal.
    \item $\Fin\oplus\mathcal{P}(\omega)$ is an ideal on $2\times\omega$ consisting of all $A\subseteq 2\times\omega$ such that $A\cap(\{0\}\times\omega)$ is finite. It is an $\bf{F_\sigma}$ non-tall P-ideal. 
    \item A summable ideal is an ideal of the form
\[
\I_{\w{r}}=\left\{A\subseteq\omega:\sum_{n\in A}r_n<\infty\right\},
\]
where $\w{r}=(r_n)$ is a sequence of positive reals such that $\sum_{n\in\omega}r_n=\infty$. Each summable ideal is an $\bf{F_\sigma}$ P-ideal. The most well-known summable ideal is
\[
\I_{1/n}=\left\{A\subseteq\omega:\sum_{n\in A}\frac{1}{n+1}<\infty\right\}.
\]
It is a tall ideal. Note that $\Fin$ is a summable ideal (given, eg. by any constant sequence $\w{r}$) and $\Fin\oplus\mathcal{P}(\omega)$ is isomorphic to the summable ideal $\{A\subseteq\omega: A\cap\{2n:n\in\omega\}\in\Fin\}$ given by the sequence:
\[
r_n=\begin{cases}
    \frac{1}{2^n} & \text{if }n\text{ is odd},\cr
    1 & \text{if }n\text{ is even}.
\end{cases}
\]
\item The ideal of asymptotic density zero is given by:
\[
\I_d= \left\{A \subseteq \omega: \frac{|A \cap (n+1)|}{n+1} \to 0 \right\}
\]
(here we use the standard set-theoretic convention and treat a number $n\in\omega$ as the set $\{0,1,\ldots,n-1\}$). It is an $\bf{F_{\sigma\delta}}$, but not $\bf{F_{\sigma}}$, tall P-ideal. 
\item $\{\emptyset\}\otimes\Fin$ is an ideal on $\omega^2$ consisting of all $A\subseteq\omega^2$ such that $A_{(n)}=\{k\in\omega: (n,k)\in A\}\in\Fin$ for all $n\in\omega$.  It is an $\bf{F_{\sigma\delta}}$, but not $\bf{F_{\sigma}}$, non-tall P-ideal. 
\end{itemize}
For more examples of ideals see for instance \cite{Farah}.
\end{example}

\subsection{Submeasures}

A function $\phi:\mathcal{P}(\omega)\to[0,\infty]$ is called a submeasure if $\phi(\emptyset)=0$, $\phi(\{n\})<\infty$, for each $n\in\omega$, and 
$$\phi(A)\leq\phi(A\cup B)\leq\phi(A)+\phi(B)$$
for all $A,B\subseteq\omega$. A submeasure $\phi$ is lower semicontinuous (lsc, in short) if $\phi(A)=\lim_{n\to\infty}\phi(A\cap n)$ for all $A\subseteq\omega$. The support of a submeasure $\phi$ is defined as $\supp(\phi)=\{n\in\omega: \phi(\{n\})\neq 0\}$.

Mazur in \cite[Lemma 1.2]{Mazur} proved that an ideal on $\omega$ is $\bf{F_\sigma}$ if and only if it is of the form:
$$\Fin(\phi)=\left\{A\subseteq\omega: \phi(A)<\infty\right\}$$
for some lower semicontinuous submeasure $\phi$ such that $\omega\notin\Fin(\phi)$ (see also \cite[Theorem 1.2.5]{Farah}). Similarly, Solecki in \cite[Theorem 3.1]{SoleckiExh} showed that an ideal on $\omega$ is an analytic P-ideal if and only if it is of the form:
$$\Exh(\phi)=\left\{A\subseteq\omega: \lim_{n\to\infty}\phi(A\setminus n)=0\right\}$$
for some lower semicontinuous submeasure $\phi$ such that $\omega\notin\Exh(\phi)$ (see also \cite[Theorem 1.2.5]{Farah}). 

A submeasure $\phi$ is non-pathological, if
$$\phi(A)=\sup\{\mu(A): \mu\text{ is a measure with }\supp(\mu)\in\Fin\text{ and }\mu(B)\leq\phi(B)\text{ for all }B\subseteq\omega\},$$
for all $A\subseteq\omega$. We say that an $\bf{F_\sigma}$ ideal (analytic P-ideal) is non-pathological, if it is of the form $\Fin(\phi)$ ($\Exh(\phi)$, respectively), for some non-pathological submeasure $\phi$.

\subsection{Ideal convergence}

Following \cite{KSW}, for an ideal $\I$ on $\omega$ and a sequence $(x_n)$ of elements of some metric space $(X,d)$, we say that $(x_n)$ is $\I$-convergent to some $x\in X$ if almost all $x_n$'s in the sense of $\I$, are arbitrarily close to $x$, that is, for every $\varepsilon>0$ we have $\{n \in \omega: d(x_n , x)> \varepsilon \} \in \I$. In such a case, we write $x_n \to_\I x$ or $\lim_{n,\I}x_n=x$. Note that for the particular ideal of asymptotic density zero, ideal convergence has been previously considered under the name of statistical convergence (see e.g. \cite{Fast}, \cite{Fridy}, \cite{Salat}, \cite{Steinhaus}).

Similarly, for a Banach space $X$, we say that a series $(\w{x}_n)$ is $\I$-convergent to some $\w{x}\in X$, if the sequence $(\w{a}_n)$ of partial sums, given by $\w{a}_n = \sum_{i=0}^n \w{x}_n$, is $\I$-convergent to $X$. In such a case we write 
\[
\w{x}=\lim_{n,\I} \w{a}_n = \sum_{n,\I} \w{x}_n.
\]
Note that it means "almost all partial sums are close to $\w{x}$" rather than "sums of almost all elements are close to $\w{x}$". In fact, even though the second statement might be more intuitive on the first thought, it does not make sense at all. 

\subsection{Schauder bases}

Given a Banach space $X$ over the real or complex field $\IK$, by a Schauder basis of $X$ we usually mean such a sequence $(\w{u}_n)\subseteq X$ that for any $\w{x}\in X$ there exists exactly one sequence of scalars $(\alpha_n)\subseteq \IK$ such that $\w{x}= \sum_{n=0}^{\infty} \alpha_n \w{u}_n$. In such a case we can associate with it a sequence of functionals $\w{u}_n^\star \colon X \to \IK $ such that $\w{u}_i^\star \colon \sum_{n=0}^{\infty} \alpha_n \w{u}_n \mapsto \alpha_i$. From Banach Space Theory, it is a standard yet nontrivial fact that all $\w{u}_n^\star$'s are continuous and linear (see eg. \cite[Theorem 1.1.3]{AK}). Throughout the paper we will be dealing with specific examples of bases and we find fixing coordinate functionals convenient, so by a Schauder basis of $X$ we mean the sequence of pairs $\hat{\w{u}}=(\w{u}_n,\w{u_n^\star})$.

A Schauder basis $\hat{\w{u}}=(\w{u}_n,\w{u_n^\star})$ of a Banach space $X$ is unconditional, if the series $\sum_{n=0}^{\infty} \w{u}_n^\star(\w{x}) \w{u}_n$ converges unconditionally, for every $\w{x}\in X$.

Throughout the whole paper by $\hat{\w{e}}=(\w{e}_n,\w{e_n^\star})$ we will denote the standard Schauder basis of $c_0$. Note that $\hat{\w{e}}$ is also a Schauder basis of each $\ell_p$, $1\leq p<\infty$. It is known that $\hat{\w{e}}$ is unconditional in $c_0$ and in each $\ell_p$, $1\leq p<\infty$.


\begin{lemma}[Folklore]
\label{lem-monotonicznosc}
Let $\hat{\w{u}}$ be an unconditional Schauder base of $X$. If $\sum_n a_n \w{u}_n\in X$ and $(b_n)\subseteq \IK$ is such that $|b_n|\leq |a_n|$, for every $n \in \omega$, then $\sum_n b_n \w{u}_n \in X$.
\end{lemma}

\begin{proof}
Since $\hat{\w{u}}$ is unconditional, by \cite[Proposition 3.1.3]{AK}, there is $K\geq 1$ such that
    \[
    \left\| \sum_{i=0}^n c_i \w{u}_i \right\|_X \leq K \left\| \sum_{i=0}^n d_i \w{u}_i \right\|_X,
\]
for all $n\in\omega$ and scalars $c_0,\ldots,c_n,d_0,\ldots,d_n$ such that $|c_i|\leq |d_i|$, for $i\leq n$.

Fix $n\in\omega$. Then for any $m>n$ 
    \[
    \left\| \sum_{i=n}^m b_n \w{u}_n \right\|_X \leq K \left\| \sum_{i=n}^m a_n \w{u}_n \right\|_X \leq  K^2 \left\| \sum_{i=n}^\infty a_n \w{u}_n \right\|_X \xrightarrow{n \to \infty} 0.
\]
Therefore $ \sum_{n=0}^\infty b_n \w{u}_n  $ is convergent.
\end{proof}



\section{Ideal Schauder bases}

\subsection{Definition}

Natural way for considering ideal version of Schauder bases is to simply demand $\w{x}=\sum_{n,\I} \alpha_n \w{e}_n$ instead of $\w{x}=\sum_{n} \alpha_n \w{e}_n$ in which case linear coordinate functionals also appear, yet their continuity is not so clear anymore (see \cite{CGK,RKS,KS,Kochanek}  for more detailed discussion of this issue). Since throughout the presented paper we find fixing coordinate functionals convenient, therefore we consider the following definition.

If $X$ is a Banach space and $\hat{\w{a}}=(\w{a}_n,\w{a}^\star_n)\subseteq X\times \IK^X$ are such that: 
\[
\w{a}^\star_n(\w{a}_m)=\delta_{nm}=\begin{cases}
    1, & \text{if }n=m,\\
    0, & \text{if }n\neq m,
\end{cases}
\]
(which represents equality $\w{a}_n=\w{a}_n +\sum_{m\neq n}0\cdot \w{a}_m$), then we denote: 
\[
S_n(\hat{\w{a}})(\w{x})=\sum_{i=0}^n \w{a}_i^\star(\w{x}) \w{a}_i,
\]
for all $\w{x}\in X$ and $n\in\omega$. For an ideal $\I$ on $\omega$, we say that $\hat{\w{a}}$ is an $\I$-Schauder basis of $X$ if
\[
\{n\in\omega: \|\w{x}-S_n(\hat{\w{a}})(\w{x})\|_X\geq\varepsilon\}\in\I,
\]
that is $\w{x}=\sum_{n,\I} \w{a}^\star_n(\w{x})\w{a}_n$, for every $\w{x}\in X$ and $\varepsilon>0$.

It is a natural question if our definition is equivalent to the original one, i.e. if the sequence $(\w{a}_n^\star (\w{x}))$ is the only possible sequence of coordinates. Now we will explain why uniqueness of coefficients is not a big deal in the context of presented paper (however, note that in general there might be an issue if some of $\w{a}_n^\star$'s would not be continuous). Our reasoning is the same as in the case of standard bases (see eg. \cite[Definition 1.1.2, Theorem 1.1.3]{AK}).

\begin{lemma}
\label{L:wspolczynniki}
Let $(\w{a}_n)$ be such a sequence in some Banach space $X$ that $\w{a}_n\notin \overline{\on{span}\{\w{a}_m:m \neq n\}}$ for every $n \in \omega$. Fix two sequences of scalars $(\alpha_n)$ and $(\beta_n)$. Then for every $l\in \omega$ such that $\alpha_l\neq \beta_l$ there exists $\varepsilon>0$ such that $\|\sum_{i=0}^k \alpha_i\w{a}_i - \sum_{i=0}^{k'} \beta_i\w{a}_i\|_X>\varepsilon$ for every $k,k'>l$.
\end{lemma}

\begin{proof}
Note that $\sum_{i=0}^k \alpha_i\w{a}_i - \sum_{i=0}^{k'} \beta_i\w{a}_i=(\alpha_l-\beta_l)\w{a}_l-\w{b}$ for some $\w{b}\in \on{span}\{\w{a}_m:m \neq l\}$. Thus
\[
\left\|\sum_{i=0}^k \alpha_i\w{a}_i - \sum_{i=0}^{k'} \beta_i\w{a}_i\right\|_X\geq \on{dist}\left((\alpha_l-\beta_l)\w{a}_l, \overline{\on{span}\{\w{a}_m:m \neq l\}}\right)>0.
\]
Note that the distance above does not depend on $k$ and $k'$. 
\end{proof}

\begin{proposition}
\label{P:wspolczynniki}
Let $\hat{\w{a}}=(\w{a}_n,\w{a}^\star_n)\subseteq X\times X^\star$ be such that $\w{a}^\star_n(\w{a}_m)=\delta_{nm}$ for all $n,m\in\omega$. Fix $\w{x}\in X$ and a sequence of scalars $(\beta_n)$. Assume there are two nontrivial ideals $\I$ and $\J$ such that:
\[
\w{x}=\sum_{n,\I} \w{a}^\star_n(\w{x}) \w{a}_n= \sum_{n,\J} \beta_n \w{a}_n.
\]
Then $\w{a}_n^\star(x)=\beta_n$ for every $n\in \omega$. 
\end{proposition}

\begin{proof}
    Note that if $n\neq m$ then the condition $\w{a}^\star_n(\w{a}_m)=0$ implies $\w{a}_m\in \on{Ker}(\w{a}^\star_n)$. By the continuity of $\w{a}_n^\star$, $\on{Ker}(\w{a}^\star_n)$ is closed, so $\overline{\on{span}\{\w{a}_m: m \neq n\}}\subseteq \on{Ker}(\w{a}^\star_n)$. Since $\w{a}_n^\star(\w{a}_n)=1$, we obtain $\w{a}_n \notin \on{ker}\w{a}^\star_n$ for every $n$. If $\beta_n\neq \w{a}_n^\star(\w{x})$ for some $\w{x}\in X$ and $n\in \omega$, then, by Lemma \ref{L:wspolczynniki}, there is some $\varepsilon>0$ such that 
    \[
    \left\|\sum_{i=0}^k \w{a}^\star_i(\w{x}) \w{a}_i- \sum_{i=0}^{k'} \beta_i \w{a}_i \right\|>\varepsilon
    \]
    for every $k,k'>n$. But since $\w{x}=\sum_{n,\I} \w{a}^\star_n(\w{x}) \w{a}_n$ and $\I\neq P(\omega)$, there exists $k>n$ such that 
    \[
     \left\|\w{x}-\sum_{i=0}^k \w{a}^\star_i(\w{x}) \w{a}_i\right\|<\frac{\varepsilon}{2},
    \]
    hence 
    \[
     \left\|\w{x}-\sum_{i=0}^{k'} \beta_i \w{a}_i\right\|>\frac{\varepsilon}{2},
    \]
    for all $k'>n$, contradicting $\w{x}= \sum_{n,\J} \beta_i \w{a}_n$.
\end{proof}

Above considerations leads us to some useful Corollaries.

\begin{corollary}
Let $X$ be a Banach space and $\I$ be a nontrivial ideal. Assume that $\hat{\w{a}}=(\w{a}_n,\w{a}^\star_n)\subseteq X\times X^\star$ is an $\I$-Schauder basis of $X$ such that $\w{a}^\star_n(\w{a}_m)=\delta_{nm}$ for all $n,m\in\omega$. If $\hat{\w{a}}$ is also a $\J$-Schauder basis of $X$, for some nontrivial ideal $\J$ (in particular, for $\J=\I$), then for every $\w{x}\in X$ there exists a unique sequence of scalars $(\alpha_n)$ ($=(\w{a}_n^\star(\w{x})))$) such that $\w{x}=\sum_{n,\J}\alpha_n \w{a}_n$.
\end{corollary}

\begin{proof}
Follows from Proposition \ref{P:wspolczynniki}.
\end{proof}

\begin{corollary}
Let $X$ be a Banach space and $\I$ be a nontrivial analytic ideal. Assume that $(\w{a}_n)\subseteq X$ is such that for every $\w{x}\in X$ there exists a unique sequence of scalars $(\alpha_{n,\w{x}})$ such that $\w{x}=\sum_{n,\I}\alpha_{n,\w{x}} \w{a}_n$. If for a nontrivial ideal $\J$ and $\w{x}\in X$ there is a sequence of scalars $(\beta_n)$ such that $\w{x}=\sum_{n,\J}\beta_n \w{a}_n$, then $(\beta_n)=(\alpha_{n,\w{x}})$.
\end{corollary}

\begin{proof}
Follows from Proposition \ref{P:wspolczynniki} and \cite[Theorem A]{RKS}.
\end{proof}


Clearly the definition of Schauder basis may be considered also in a normed vector spaces, and we will do so in Section \href{https://arxiv.org/abs/2508.21235v1}{10}.

\subsection{Critical ideal}

It is known (see \cite[Proof of Theorem B]{RKS}) that for each Banach space $X$ and each $\hat{\w{a}}=(\w{a}_n,\w{a}^\star_n)\subseteq X\times \IK^X$ there is an ideal $\CR(X,\hat{\w{a}})$ on $\omega$ such that the following are equivalent for every ideal $\I$ on $\omega$:
\begin{itemize}
    \item $\hat{\w{a}}$ is an $\I$-Schauder basis of $X$,
    \item $\CR(X,\hat{\w{a}})\subseteq\I$.
\end{itemize}
We say that $\hat{\w{a}}$ is an ideal Schauder basis of $X$ if $\CR(X,\hat{\w{a}})\neq\mathcal{P}(\omega)$. Obviously, if $\CR(X,\hat{\w{a}})=\Fin$, then $\hat{\w{a}}$ is a Schauder basis of $X$.

\subsection{Simple ideal Schauder bases}

In this paper we are mainly interested in ideal Schauder bases of special form:

\begin{definition}
\label{def-simple}
Let $X$ be a vector space, $\hat{\w{a}}=(\w{a}_n,\w{a}^\star_n)\subseteq X\times X^\star$ and $\hat{\w{u}}=(\w{u}_n,\w{u}^\star_n)\subseteq X\times X^\star$. We say that $\hat{\w{a}}$ is simple over $\hat{\w{u}}$ if there are $D\subseteq\omega$, an interval-to-one $h:D\to\omega$ (i.e., $h^{-1}[\{k\}]$ is either empty or an interval contained in $D$, for every $k\in\omega$) and $(\w{b}_n)_{n\in h[D]}\subseteq X\setminus\{0\}$ such that:
\[
S_n(\hat{\w{u}})(\w{x})-S_n(\hat{\w{a}})(\w{x})=
\begin{cases}
    0, & \text{if }n\in\omega\setminus D,\\
    \w{u}_{h(n)}^\star(\w{x})\w{b}_{h(n)}, & \text{if }n\in D,
\end{cases}
\]
for all $n\in \omega$ and $\w{x}\in X$. In particular, if $X$ is a Banach space and $\hat{\w{u}}$ is a $\Fin$-Schauder basis of $X$, then for a given ideal $\I$, $(S_n(\hat{\w{a}})(\w{x}))$ is $\I$-convergent to $\w{x}$ if and only if $(\w{u}_{h(n)}^\star(\w{x})\w{b}_{h(n)})_{n\in D}$ is $\I\restriction D$-convergent to zero in $X$.
\end{definition}

However the above definition may be considered technical, we decided to restrict our considerations to such ideal Schauder bases for two reasons: all our examples are of this form and this special form enabled us to classify all such ideal Schauder bases in standard Banach spaces.

Next result, in some sense, allows us to consider only a special class of Banach spaces.

\begin{theorem}
If $\hat{\w{c}}$ is simple over some normed $\Fin$-Schauder basis $\hat{\w{u}}$ (that is, $\|\w{u}_n\|=1$ for each $n$) of some Banach space $Y$ over $\mathbb{R}$, then $\CR(Y,\hat{\w{c}})=\CR(X,\hat{\w{a}})$ for some Banach space $\ell_1\subseteq X\subseteq c_0$ (equipped in coordinate-wise addition and multiplication by scalars) and some $\hat{\w{a}}$ simple over $\hat{\w{e}}$. Moreover:
\begin{itemize}
    \item[(i)] $\hat{\w{u}}$ is unconditional in $Y$ if and only if $\hat{\w{e}}$ is unconditional in $X$,
    \item[(ii)] the topology on $X$ is stronger than the topology inherited from $c_0$,
    \item[(iii)] the topology on $X$ is weaker than $\ell_1$ topology on $X$ (understood as topology generated by sets of the form $B((x_n),r)=\{(y_n)\in X : \sum_{n\in\omega} |x_n-y_n|<r\}$).
\end{itemize}
\end{theorem}

\begin{proof}
If $Y$ is an arbitrary Banach space over $\mathbb{R}$ and $\hat{\w{u}}$ is its $\Fin$-Schauder basis, then consider: 
\[
X=\left\{\w{x}\in\mathbb{R}^\omega: \sum_{n\in\omega}x_n\w{u}_n\in Y\right\}
\]
normed by: 
\[
\|\w{x}\|_X=\left\|\sum_{n\in\omega}x_n\w{u}_n\right\|_Y,
\]
for all $\w{x}\in X$. Then $X$ is a Banach space and $\ell_1\subseteq X\subseteq c_0$.

If $\hat{\w{c}}$ is simple over $\hat{\w{u}}$, then $\CR(Y,\hat{\w{c}})=\CR(X,\hat{\w{a}})$, where $\hat{\w{a}}$ is simple over $\hat{\w{e}}$ with the same witnessing $D$, $h$ and $(\w{b}_n)$ as for $\hat{\w{c}}$ and $\hat{\w{u}}$.

Item (i) follows from the fact that $X$ and $Y$ are isometric (by the isometry that moves $\hat{\w{e}}$ to $\hat{\w{u}}$). To show item (ii), observe that for any $\w{x}_n, \w{x}\in X$ and any $k\in \IN$ we have:
\begin{align*}
    |\w{u}^\star_{k+1}(\w{x}_n-\w{x})| & =  \|\w{u}^\star_{k+1}(\w{x}_n-\w{x})\w{u}_{k+1}\|_X = \|S_{k+1}(\hat{\w{u}}) (\w{x}_{n}-\w{x})- S_{k}(\hat{\w{u}}) (\w{x}_{n}-\w{x})\|_X=\cr
    & = \|(S_{k+1}-S_k)(\hat{\w{u}}) (\w{x}_{n}-\w{x})\|_X \leq 2K \|\w{x}_{n}-\w{x}\|_X,
\end{align*}
where $K$ is the basis constant of $\hat{\w{u}}$ (see \cite[Proposition 1.1.4 and Definition 1.1.5]{AK}). The final claim is clear.
\end{proof}




\section{Some lemmas}

\begin{lemma}
\label{general-lem1}
Let $\hat{\w{a}}$ be simple over a $\Fin$-Schauder basis $\hat{\w{u}}$ of a Banach space $X$, and let $D$, $h$ and $(\w{b}_n)_{n\in\omega}$ be as in Definition \ref{def-simple}. Then $\CR(X,\hat{\w{a}})$ is the ideal on $\omega$ generated by $\Fin$ and all sets of the form
    \[
    A_{\w{x}}=\left\{n\in D: |\w{u}^\star_{h(n)}(\w{x})|\geq\frac{1}{\|\w{b}_{h(n)}\|_X}\right\},
    \]
for $\w{x}\in X$.
\end{lemma}

\begin{proof}
Let $\hat{\J}(X,\hat{\w{u}},\hat{\w{a}})$ be the ideal on $\omega$ generated by $\Fin$ and all sets of the form $A_{\w{x}}$, for $\w{x}\in X$. We need to show that the following are equivalent for any ideal $\I$ on $\omega$:
\begin{itemize}
    \item[(a)] $\hat{\w{a}}$ is an $\I$-Schauder base of $X$;
    \item[(b)] $\hat{\J}(X,\hat{\w{u}},\hat{\w{a}})\subseteq\I$.
\end{itemize}

(a)$\implies$(b): Assume that $\hat{\J}(X,\hat{\w{u}},\hat{\w{a}})\not\subseteq\I$. Then there exists an $A\in\hat{\J}(X,\hat{\w{u}},\hat{\w{a}})\setminus\I$. Hence, there are $F\in\Fin$, $k\in\omega$ and $\w{x}_0,\ldots,\w{x}_k\in X$ such that $A\subseteq F\cup A_{\w{x}_0}\cup\ldots\cup A_{\w{x}_k}$. Since $A\notin\I$, we can find $j\leq k$ such that $A_{\w{x}_j}\notin\I$. Observe that $(\w{u}_{h(n)}^\star(\w{x}_j)\w{b}_{h(n)})_{n\in D}$ is not $\I\restriction D$-convergent to zero, as:
\[
\{n\in D: \|\w{u}_{h(n)}^\star(\w{x}_j)\w{b}_{h(n)}\|_X\geq 1\}=A_{\w{x}_j}\notin\I.
\]
Thus, $\hat{\w{a}}$ is not an $\I$-Schauder basis of $X$. 

(b)$\implies$(a): Assume that $\hat{\J}(X,\hat{\w{u}},\hat{\w{a}})\subseteq\I$ and fix any $\w{x}\in X$ and $\varepsilon>0$. Then $\frac{\w{x}}{\varepsilon}\in X$, so $A_{\frac{\w{x}}{\varepsilon}}\in\hat{\J}(X,\hat{\w{u}},\hat{\w{a}})\subseteq\I$. To finish the proof, note that:
\[
\{n\in D: \|\w{u}_{h(n)}^\star(\w{x})\w{b}_{h(n)}\|_X\geq \varepsilon\}=\left\{n\in D: \left|\w{u}_{h(n)}^\star\left(\frac{\w{x}}{\varepsilon}\right)\right|\geq \frac{1}{\|\w{b}_{h(n)}\|_X}\right\}=A_{\frac{\w{x}}{\varepsilon}}.
\]
\end{proof}

\begin{remark}
\label{remark-general-lem1}
Note that the above Lemma holds in all normed spaces. We will use this fact in Section \href{https://arxiv.org/abs/2508.21235v1}{10}.
\end{remark}

\begin{lemma}
\label{general-lem2}
Let $\hat{\w{u}}$ be an unconditional $\Fin$-Schauder basis of some Banach space $X$. Assume that $\hat{\w{a}}$ is simple over $\hat{\w{u}}$ and $D$, $h$ and $(\w{b}_n)_{n\in\omega}$ are as in Definition \ref{def-simple}. Then:
\[
\CR(X,\hat{\w{a}})=\left\{A\subseteq\omega: A\setminus D\in\Fin\text{ and }\sum_{k\in h[A\cap D]}\frac{\w{u}_{k}}{\|\w{b}_{k}\|_X}\in X\right\}.
\]
\end{lemma}

\begin{proof}
By Lemma \ref{general-lem1}, $\CR(X,\hat{\w{a}})$ is the ideal on $\omega$ generated by $\Fin$ and all sets of the form
    \[
    A_{\w{x}}=\left\{n\in D: |\w{u}^\star_{h(n)}(\w{x})|\geq\frac{1}{\|\w{b}_{h(n)}\|_X}\right\},
    \]
for $\w{x}\in X$. Let 
\[
\J(X,\hat{\w{u}},\hat{\w{a}})=\left\{A\subseteq\omega: A\setminus D\in\Fin\text{ and }\sum_{k\in h[A\cap D]}\frac{\w{u}_{k}}{\|\w{b}_{k}\|_X}\in X\right\}.
\]

The inclusion $\J(X,\hat{\w{u}},\hat{\w{a}})\subseteq\CR(X,\hat{\w{a}})$ is obvious, as $A\in \J(X,\hat{\w{u}},\hat{\w{a}})$ means that $A\setminus D\in\Fin\subseteq\CR(X,\hat{\w{a}})$ and $\w{x}=\sum_{k\in h[A\cap D]}\frac{\w{u}_{k}}{\|\w{b}_{k}\|_X}\in X$, which gives us $A\cap D\subseteq A_{\w{x}}\in \CR(X,\hat{\w{a}})$.

Now we show $\J(X,\hat{\w{u}},\hat{\w{a}})\supseteq\CR(X,\hat{\w{a}})$. It is obvious that $\CR(X,\hat{\w{a}})\restriction\omega\setminus D=\Fin\restriction\omega\setminus D\subseteq \J(X,\hat{\w{u}},\hat{\w{a}})$. 

Let $A\in \CR(X,\hat{\w{a}})\restriction D$. Then there are $l\in\omega$ and $\w{x}_0,\ldots,\w{x}_l\in X$ such that $A\subseteq A_{\w{x}_0}\cup\ldots\cup A_{\w{x}_l}$. Define $B_0=h[A_{\w{x}_0}]$ and $B_i=h[A_{\w{x}_i}]\setminus\bigcup_{j<i}h[A_{\w{x}_j}]$ for $i=1,\ldots,l$. Note that for each $i\leq l$, since $\w{x}_i\in X$, we have:
\[
\w{y}_i=\sum_{k\in B_i\cap h[A\cap D]}\w{u}_{k}^\star(\w{x}_i)\w{u}_{k}\in X
\]
(by Lemma \ref{lem-monotonicznosc}). Hence, also $\w{z}_i=\sum_{k\in B_i\cap h[A\cap D]}\frac{\w{u}_{k}}{\|\w{b}_{k}\|}\in X$ (again, by Lemma \ref{lem-monotonicznosc}). Finally, since the sets $B_i$ are pairwise disjoint and $h[A\cap D]=\bigcup_{i\leq l}h[A\cap D]\cap B_i$, we get $\sum_{k\in h[A\cap D]}\frac{\w{u}_{k}}{\|\w{b}_{k}\|}\in X$.
\end{proof}

Next example shows that in Lemma \ref{general-lem2} the assumption about unconditionality of $\hat{\w{u}}$ cannot be dropped. 

\begin{example}
Consider $\hat{\w{u}}$ given by $\w{u}_n = \sum_{i=0}^n \w{e}_n $ and $\w{u}_n^\star=\w{e}_n^\star - \w{e}_{n+1}^\star$, for all $n\in\omega$, which is a non-unconditional $\Fin$-Schauder basis of $c_0$ (see \cite[Example 3.1.2]{AK}).

Let $\hat{\w{a}}$ be given by $\w{a}_n = \sum_{i=0}^n \w{u}_n $ and $\w{a}_n^\star=\w{u}_n^\star - \w{u}_{n+1}^\star$, for all $n\in\omega$.

Observe that:
\[
S_n(\hat{\w{a}})(\w{x})=\sum_{i=0}^n (\w{u}_i^\star-\w{u}_{i+1}^\star) (\w{x})\sum_{k=0}^i \w{u}_k = \sum_{i=0}^n \w{u}_i^\star (\w{x}) \w{u}_i - \w{u}_{n+1}^\star(\w{x}) \sum_{i=0}^n \w{u}_i=S_n(\hat{\w{u}})(\w{x})-\w{u}_{n+1}^\star(\w{x}) \sum_{i=0}^n \w{u}_i,
\]
for any $\w{x}\in c_0$. Thus, 
\[
S_n(\hat{\w{u}})(\w{x})-S_n(\hat{\w{a}})(\w{x})=\w{u}_{n+1}^\star(\w{x}) \sum_{i=0}^n \w{u}_i
\]
and $\hat{\w{a}}$ is simple over $\hat{\w{u}}$ (as witnessed by $D=\omega$, $h(n)=n+1$ and $\w{b}_{n+1}=\sum_{i=0}^n \w{u}_i$).

Note that: 
\[
\omega\notin\left\{A\subseteq\omega: A\setminus D\in\Fin\text{ and }\sum_{k\in h[A\cap D]}\frac{\w{u}_{k}}{\|\w{b}_{k}\|_{c_0}}\in c_0\right\}.
\]
Indeed, $\w{x}=\sum_{k\in h[\omega]}\frac{\w{u}_{k}}{\|\w{b}_k\|_{c_0}}=\sum_{n\in\omega}\frac{\w{u}_{n+1}}{n+1}\notin c_0$, since $\w{e}_0(\w{x})=\sum_{n\in\omega}\frac{1}{n+1}=\infty$.

We will show that $\omega\in\CR(c_0,\hat{\w{a}})$. Consider $\w{x}=\sum_{n\in\omega}\frac{\w{e}_{2n}}{n+1}\in c_0$ and $\w{y}=\sum_{n\in\omega}\frac{\w{e}_{2n+1}}{n+1}\in c_0$. We have:
\[
A_{\w{x}}=\left\{n\in\omega: |\w{u}^\star_{n+1}(\w{x})|\geq\frac{1}{\|\w{b}_{n+1}\|_{c_0}}\right\}=\left\{n\in\omega: |\w{e}_{n+1}^\star(\w{x}) - \w{e}_{n+2}^\star(\w{x})|\geq\frac{1}{n+1}\right\}.
\]
Note that if $n\in\{2k+1:k\in\omega\}$, i.e., $n=2m+1$ for some $m\in\omega$, then 
\[
|\w{e}_{n+1}^\star(\w{x}) - \w{e}_{n+2}^\star(\w{x})|=|\w{e}_{2m+2}^\star(\w{x}) - \w{e}_{2m+3}^\star(\w{x})|=\w{e}_{2m+2}^\star(\w{x})=\frac{1}{m+2}\geq \frac{1}{n+1}.
\]
This shows that $\{2k+1:k\in\omega\}\subseteq A_{\w{x}}$.

Similarly one can show that $\{2k:k\in\omega\}\subseteq A_{\w{y}}$. Hence, applying Lemma \ref{general-lem1}, $\omega=\{2k:k\in\omega\}\cup\{2k+1:k\in\omega\}\in\CR(c_0,\hat{\w{a}})$.
\end{example}

\begin{proposition}
\label{general-P-ideal}
If $\hat{\w{u}}$ is an unconditional $\Fin$-Schauder basis of some Banach space $X$ and $\hat{\w{a}}$ is simple over $\hat{\w{u}}$, then $\CR(X,\hat{\w{a}})$ is a P-ideal.
\end{proposition}

\begin{proof}
From Lemma \ref{general-lem2} we know that:
\[
\CR(X,\hat{\w{a}})=\left\{A\subseteq\omega: A\setminus D\in\Fin\text{ and }\sum_{k\in h[A\cap D]}\frac{\w{u}_{k}}{\|\w{b}_{k}\|_X}\in X\right\},
\]
where $D$, $h$ and $(\w{b}_n)_{n\in\omega}$ are as in Definition \ref{def-simple}. 

Let $(A_k)\subseteq \CR(X,\hat{\w{a}})$. For each $k\in\omega$, since $\sum_{i\in h[A_k\cap D]}\frac{\w{u}_{i}}{\|\w{b}_{i}\|_X}\in X$ and $h$ is interval-to-one, there is $m_k\in\omega$ such that $\left\|\sum_{i\in h[(A_k\cap D)\setminus m_k]}\frac{\w{u}_{i}}{\|\w{b}_{i}\|_X}\right\|_X<\frac{1}{2^k}$. Define $A=\bigcup_{k\in\omega}((A_k\cap D)\setminus m_k)$. Then $A\subseteq D$ and $A_k\setminus A\in\Fin$, for each $k$. We need to show that $A\in\CR(X,\hat{\w{a}})$, i.e., $\sum_{i\in h[A]}\frac{\w{u}_{i}}{\|\w{b}_{i}\|_X}\in X$. 

Given any $\varepsilon>0$, find $j\in\omega\setminus\{0\}$ such that $\frac{1}{2^{j-1}}<\varepsilon$ and $m\in\omega$ such that 
\[
\sum_{k\leq j}\left\|\sum_{i\in h[A_k\cap D]\setminus m}\frac{\w{u}_{i}}{\|\w{b}_{i}\|_X}\right\|_X<\frac{1}{2^{j}}.
\]
Observe that
\begin{align*}
    \left\|\sum_{i\in h[A]\setminus m}\frac{\w{u}_{i}}{\|\w{b}_{i}\|_X}\right\|_X & \leq\sum_{k\leq j}\left\|\sum_{i\in h[A\cap A_k]\setminus m}\frac{\w{u}_{i}}{\|\w{b}_{i}\|_X}\right\|_X+\sum_{k>j}\left\|\sum_{i\in h[A\cap A_k]}\frac{\w{u}_{i}}{\|\w{b}_{i}\|_X}\right\|_X< \cr
    & <\frac{1}{2^{j}}+\sum_{k>j}\frac{1}{2^k}=\frac{1}{2^{j-1}}<\varepsilon.
\end{align*}
\end{proof}



\section{Banach spaces with a Schauder basis}


\begin{theorem}
\label{thm-distinguishing}
Let $X$ be any Banach space with a $\Fin$-Schauder basis $\hat{\w{u}}$. For any infinite $E\subseteq\omega$ there is $\hat{\w{a}}\subseteq X\times X^\star$ simple over $\hat{\w{u}}$ such that the following are equivalent for any ideal $\I$:
\begin{itemize}
    \item[(a)] $\hat{\w{a}}$ is an $\I$-Schauder basis of $X$;
    \item[(b)] $E\in\I$;
    \item[(c)] $\I_E\subseteq\I$, where $\I_E$ is the ideal generated by $\Fin\cup\{E\}$.
\end{itemize}
In particular, $\CR(X,\hat{\w{a}})=\I_E$.
\end{theorem}

\begin{proof}
Fix a $\Fin$-Schauder basis $\hat{\w{u}}$ of $X$. Without loss of generality we may assume that $\hat{\w{u}}$ is normed, that is, $\|\w{u}_n\|_X=1$ for all $n\in\omega$. Let $(I_k)$ be a sequence of pairwise disjoint intervals in $\omega$ such that $E=\bigcup_k I_k$ and $\max I_k+1<\min I_{k+1}$. Denote $b_k=\min I_k$ and let 
\[
d_k=\min\left\{\left\|\sum_{i=j}^{\max I_k+1} \w{u}_{i}\right\|_X: b_k<j\leq \max I_k+1\right\},
\]
for all $k\in\omega$.

Define 
\[
    \w{a}_n =
    \begin{cases*}
    \w{u}_{n}+\frac{2^n}{d_k}\cdot\sum_{i=n+1}^{\max I_k+1} \w{u}_{i}     & if $n=b_k$ for some $k$, \\
    \w{u}_n & otherwise,  \\
    \end{cases*}
  \]
and
\[
    \w{a}_n^\star =
    \begin{cases*}
    \w{u}_n^\star & if  $n-1\notin E$  \\
    \w{u}_{n}^\star - \frac{2^{b_k}}{d_k}\cdot \w{u}_{b_k}^\star & if $n-1\in I_k$ for some $k$.
    \end{cases*}
  \]

Clearly, $\w{a}^\star_n(\w{a}_m)=\delta_{nm}$ for all $n,m\in\omega$.

The equivalence of items (b) and (c) is obvious. We will show that items (a) and (b) are equivalent.

Observe that:
\[
S_n(\hat{\w{u}})(\w{x})-S_n(\hat{\w{a}})(\w{x})=\begin{cases*}
    0, & if  $n\notin E$  \\
    \w{u}^\star_{b_k}(\w{x})\frac{2^{b_k}}{d_k}\sum_{i=j}^{\max I_k+1}\w{u}_i, & if $n=b_k+j$ for some $k\in\omega$ and $0\leq j<|I_k|$.
    \end{cases*}
\]
Hence, if $E\in\I$, then $\hat{\w{a}}$ is an $\I$-Schauder basis of $X$. 

We need to show that if $E\notin\I$, then $\hat{\w{a}}$ is not an $\I$-Schauder basis of $X$. Define $\w{x}=\sum_{n\in\omega}\frac{1}{2^n}\w{u}_n$ (note that $\w{x}\in X$, since $\|\sum_{n>k}\frac{1}{2^n}\w{u}_n\|_X\leq\frac{1}{2^k}$ for all $k\in\omega$). Then 
\[
\{n\in\omega: \|S_n(\hat{\w{u}})(\w{x})-S_n(\hat{\w{a}})(\w{x})\|_X\geq 1\}\supseteq E\notin\I,
\]
since given any $n\in E$, find $k\in\omega$ and $0\leq j<|I_k|$ such that $n=b_k+j\in I_k$, and observe that:
\[
\|S_n(\hat{\w{u}})(\w{x})-S_n(\hat{\w{a}})(\w{x})\|_X=\left\|\w{u}^\star_{b_k}(\w{x})\frac{2^{b_k}}{d_k}\sum_{i=j}^{\max I_k+1}\w{u}_i\right\|_X\geq |2^{b_k}\w{u}^\star_{b_k}(\w{x})|=1.
\]
\end{proof}

\begin{thebibliography}{99}
\bibitem[1]{AK} F.~Albiac, N.J.~Kalton, \emph{Topics in {B}anach space theory. Second edition} With a foreword by Gilles Godefory, Graduate Texts in Mathematics, 233. Springer, [Cham], 2016.
\bibitem[5]{CGK} J.~Connor, M.~Ganichev, and V.~Kadets, \emph{A characterization of Banach spaces with separable duals via weak statistical convergence}, J. Math. Anal. Appl. \textbf{244} (2000), 251--261.
\bibitem[6]{Debs} G.~Debs, J.~Saint Raymond, \emph{Filter descriptive classes of Borel functions}, Fundamenta Mathematicae 204 (2009), 189--213.
\bibitem[7]{RKS} N.~de~Rancourt, T.~Kania and J.~Swaczyna, Continuity of coordinate functionals of filter bases in Banach spaces. \emph{J. of Funct. Anal.}. \textbf{284} (2023), no. 9, 109869.
\bibitem[8]{Farah} I.~Farah, \emph{Analytic quotients. Theory of lifting for quotients over analytic ideals on integers},
Mem. Amer. Math. Soc. 148 (2000).
\bibitem[9]{Fast} H. Fast, \emph{Sur la convergence statistique}, Colloq. Math. 2 (1951), 41--44.
\bibitem[10]{Fridy} J. A. Fridy, \emph{On statistical convergence}, Analysis 5 (1985), 301--313.
\bibitem[14]{KS} T.~Kania and J.~Swaczyna, Large cardinals and continuity of coordinate functionals of filter bases in Banach spaces. \emph{Bull. Lond. Math. Soc}. \textbf{53} (2021), no. 1, 231--239.
\bibitem[15]{Kochanek} T.~Kochanek, \emph{$\mathcal F$-bases with brackets and with individual brackets in Banach spaces}, Studia Math. 211 (2012), 259--268.
\bibitem[16]{KSW} P. Kostyrko, T. \v{S}alat, W. Wilczy\'nski, \emph{I-convergence}, Real Anal. Exchange 26 (2000/01), 669-685.
\bibitem[17]{Mazur} K. Mazur, \emph{{$F_\sigma$}-ideals and {$\omega_1\omega_1^*$}-gaps in the {B}oolean algebras {$P(\omega)/I$}}, Fund. Math. 138 (1991), 103--111.
\bibitem[18]{Salat} T. \v{S}al\'{a}t,
\emph{On statistically convergent sequences of real numbers},
Math. Slovaca 30 (1980), 139--150.
\bibitem[19]{Solecki} S. Solecki, \emph{Analytic ideals}, The Bulletin of Symbolic Logic 2, (1996), 339--348.
\bibitem[20]{SoleckiExh} S. Solecki, \emph{Analytic ideals and their applications}, Ann. Pure Appl. Logic 99 (1999), 51--72.
\bibitem[21]{Steinhaus} H. Steinhaus,
\emph{Sur la convergence ordinaire et la convergence asymptotique},
Colloq. Math. 2 (1951), 73--74.
\end{thebibliography}
\end{document}
